% Bagian: incompleteness
% Bab: theories-computability
% Subbagian: first-incompleteness

\documentclass[../../../include/open-logic-section]{subfiles}

\begin{document}

\olfileid[id]{inc}{tcp}{inc}

\olsection{$\Th{Q}$ Tidak Memiliki Perluasan Lengkap dan Konsisten yang
  !!^{axiomatizable}}

\begin{thm}
\ollabel{thm:first-incompleteness}
Tidak ada perluasan $\Th{Q}$ yang lengkap, konsisten, dan
!!{axiomatizable}.
\end{thm}

\begin{proof}
Kita telah mengetahui bahwa tidak ada perluasan $\Th{Q}$ yang konsisten dan
dapat diputuskan. Namun, jika $\Th{T}$ lengkap dan !!{axiomatized} secara
komputabel, maka teori itu dapat diputuskan.
\end{proof}

\begin{explain}
Teorema ini tidak terlalu jauh dari rumusan asli G\"odel pada 1931 tentang
Teorema Ketaklengkapan Pertama. Selain istilah yang lebih modern, perbedaan
utamanya ialah sebagai berikut: G\"odel menggunakan
``$\omega$-konsisten'' alih-alih ``konsisten''; dan ia belum dapat mengatakan
``!!{axiomatizable}'' dalam keumuman penuh karena gagasan formal
komputabilitas belum tersedia. (Model-model formal komputabilitas dikembangkan
sepanjang dasawarsa berikutnya, termasuk oleh G\"odel, dan sebagian besar
untuk dapat mencirikan jenis teori yang terkena fenomena G\"odel.)

Teorema tersebut menyatakan bahwa kita tidak dapat memperoleh semuanya
sekaligus, yakni kelengkapan, konsistensi, dan !!{axiomatizability}. Namun,
jika salah satunya dilepaskan, dua sifat lainnya dapat diperoleh: $\Th{Q}$
konsisten dan diaksiomatisasi secara komputabel, tetapi tidak lengkap; teori
tak konsisten lengkap dan diaksiomatisasi secara komputabel (misalnya oleh
$\{ \eq/[0][0] \}$), tetapi tidak konsisten; dan himpunan kalimat aritmetika
yang benar lengkap dan konsisten, tetapi tidak diaksiomatisasi secara
komputabel.
\end{explain}
\end{document}
